\documentclass[11pt]{article}
\usepackage{titlesec}
\titleformat{\section}[hang]{\normalfont\scshape}{\thesection.}{1em}{}
\titleformat{\subsection}[hang]{\normalfont\scshape}{\thesubsection.}{1em}{}
\usepackage[T1]{fontenc}
\usepackage[dvipsnames]{xcolor}
\usepackage[margin=2cm]{geometry}
\usepackage{amssymb} 
\usepackage{amsmath}
\usepackage{graphicx}
\usepackage{float}
\usepackage[normalem]{ulem} 
\usepackage{enumitem} 
\setlist[enumerate]{itemsep=0mm}
\usepackage{xcolor}
\setenumerate{label=(\arabic*)}
\usepackage[utf8]{inputenc}
\usepackage[T1]{fontenc}
\usepackage{babel}
\usepackage{mathtools}
\usepackage{amsthm}
\usepackage{thmtools}
\usepackage{etoolbox}
\usepackage{fancybox}
\usepackage{framed}
\usepackage{tcolorbox}
\usepackage{xcolor} 
% open 
\newcommand{\open}[0]{\mathcal{O}}
\newcommand{\topo}[0]{\mathcal{T}}
\newcommand{\hood}[0]{\mathcal{U}}
\newcommand{\base}[0]{\mathcal{B}} 
% example environmentq
\theoremstyle{definition} 
\newtheorem{exmp}{Example}[section]

% linear operators
\newcommand{\lop}[2]{\mathcal{L}(#1, #2)}

% question environment
\theoremstyle{definition}
\newtheorem{question}{Question}

%rd 
\newcommand{\rd}[0]{\mathbb{R}^d}
\newcommand{\R}[0]{\mathbb{R}}
\newcommand{\N}[0]{\mathbb{N}} 
% let E in R be measurable 
\newcommand{\EinR}[0]{Let $E \subseteq \R$ be measurable}

%integral 
\newcommand{\idx}[2]{\int_{#1}^{#2}}

% weak convergence 
\newcommand{\warrow}[0]{\rightharpoonup}
\newcommand{\fcvw}[0]{ \{f_n \} \warrow f \text{ in } L^p(E)} 

% Probability 
\DeclareRobustCommand{\bbone}{\text{\usefont{U}{bbold}{m}{n}1}}
\newcommand{\Var}[1]{\mathrm{Var[#1]}}			% variance
\newcommand{\EX}[1]{\mathbb{E}\mathrm{[#1]}}	 % expected value 
\newcommand{\seq}[1]{\{ #1_n	\}_{n \in \bb{N}}} % sequence of events
\newcommand{\pspace}[0]{( \Omega, F, P)}		% probability space
\newcommand{\msp}[0]{( \Omega, F)}		% measurable space
	
% Exercise environment 
\newenvironment{myleftbar}{%
\def\FrameCommand{\hspace{0.6em}\vrule width 2pt\hspace{0.6em}}%
\MakeFramed{\advance\hsize-\width \FrameRestore}}%
{\endMakeFramed}
\declaretheoremstyle[
spaceabove=6pt,
spacebelow=6pt
headfont=\normalfont\bfseries,
headpunct={} ,
headformat={\cornersize*{2pt}\ovalbox{\NAME~\NUMBER\ifstrequal{\NOTE}{}{\relax}{\NOTE}:}},
bodyfont=\normalfont,
]{exobreak}

\declaretheorem[style=exobreak, name=Exercise,%
postheadhook=\leavevmode\myleftbar, %
prefoothook = \endmyleftbar]{exo}

% Solution environment 
\newenvironment{mysolbar}{%
\def\FrameCommand{\hspace{0.6em}\vrule width 2pt\hspace{0.6em}}%
\MakeFramed{\advance\hsize-\width \FrameRestore}}%
{\endMakeFramed}
\declaretheoremstyle[
spaceabove=6pt,
spacebelow=6pt
headfont=\normalfont\bfseries,
headpunct={} ,
headformat={\cornersize*{2pt}\ovalbox{\NAME~\NUMBER\ifstrequal{\NOTE}{}{\relax}{\NOTE}:}},
bodyfont=\normalfont,
]{solbreak}

\declaretheorem[style=solbreak, name=Solution,%
postheadhook=\leavevmode\mysolbar, %
prefoothook = \endmysolbar]{sol}

% HEADERS
\usepackage{fancyhdr}
 
\pagestyle{fancy}
\fancyhf{}
\fancyhead[LE,RO]{Page \thepage}
\fancyhead[RE,LO]{Math 358: Advanced Calculus}
\fancyhead[CE,CO]{Winter 2026 -- Final Summary}
\fancyfoot[LE,RO]{}

% Definitions
\newcommand{\dfn}[1]{\underline{\textbf{#1}}}
\newcommand{\im}[1]{\textbf{\textcolor{red}{#1}}}

% lower integral
\usepackage{accents}

\newcommand{\ubar}[1]{\underaccent{\bar}{#1}}
\def\avint{\mathop{\,\rlap{-}\!\!\int}\nolimits} 

% custom commands 
\newcommand{\bb}[1]{\mathbb{#1}}
\newcommand{\vc}[1]{\mathbf{#1}}
\newcommand{\step}[1]{\textbf{#1}\textbf{. Step:}}
\newcommand{\pdv}[2]{\frac{\partial #1}{\partial #2}}
\newcommand{\sets}[2]{ \left\{ #1\ |\ #2 \right\}}
\DeclareMathOperator{\Tr}{Tr}

% Custom operators for Advanced cal
\DeclareMathOperator{\Vol}{Vol}
\DeclareMathOperator{\Content}{Content}
\DeclareMathOperator{\Area}{Area}
\DeclareMathOperator{\curl}{curl}
\DeclareMathOperator{\dive}{div}

% Proofs
\newcommand{\claim}[1]{\textbf{#1}\textbf{. Claim:}}

	% iff proofs
	\newcommand{\rhs}[0]{(\Rightarrow )}
	\newcommand{\lhs}[0]{(\Leftarrow )}

% sequence of functions
\newcommand{\funcseqx}{(f_n(x))_{n \in \bb{N}}}
\newcommand{\funcseq}{(f_n)_{n \in \bb{N}}}

% measurable sets 
\newcommand{\measurable}{f^{-1}([-\infty, c[)} 

% heat equation 
\newcommand{\pbdry}[2]{C^{(#1, #2)} (\Omega_T) \cap C (\overline{\Omega_T})}
\DeclareMathOperator\erf{erf}
\newcommand{\mbf}[1]{\mathbf{#1}}

% Laplace Equation 
\newcommand{\lapbdry}[1]{C^{#1} (\Omega) \cap C (\overline{\Omega})}


% math environments 
\usepackage[utf8]{inputenc}
\newtheorem{theorem}{\textcolor{blue}{Theorem}}
\newtheorem{corollary}{Corollary}
\newtheorem{lemma}[theorem]{Lemma}
\theoremstyle{definition}
\newtheorem{definition}{\textcolor{OliveGreen}{Definition}}
\newtheorem{prop}{\textcolor{red}{Proposition}}
\newtheorem{ex}{\textcolor{Maroon}{Example}}
\theoremstyle{remark}
\newtheorem*{remark}{Remark}

% cookbook proofs 
\newcommand{\cb}[3]{\underline{(#1 #2): #3:}}

\usepackage{tcolorbox}
\tcbuselibrary{theorems}

% theorems 
\newtcbtheorem[number within=section]{mytheo}{Theorem}%
{colback=blue!5,colframe=blue!35!black,fonttitle=\bfseries}{th}

% definitions 
\newtcbtheorem[number within=section]{defn}{Definition}%
{colback=black!5,colframe=black!35!black,fonttitle=\bfseries}{th}

% axioms
\newtcbtheorem[number within=section]{ax}{Axioms}%
{colback=OliveGreen!5,colframe=black!35!OliveGreen,fonttitle=\bfseries}{th}


% important examples
\newtcbtheorem[number within=section]{examp}{Example}%
{colback=Mahogany!5,colframe=black!35!Mahogany,fonttitle=\bfseries}{th}

% upper and lower riemann integrals
\newcommand{\upRiemannint}[2]{
  \overline{\int_{#1}^{#2}}
}
\newcommand{\loRiemannint}[2]{
  \underline{\int_{#1}^{#2}}
}

\usepackage{hyperref}
\hypersetup{
    colorlinks,
    citecolor=black,
    filecolor=black,
    linkcolor=black,
    urlcolor=black
}

\begin{document}

\begin{center}
	\textbf{Math 358: Advanced Calculus Summary} \\
	\textbf{Final Exam Date: 22 April 2026 6.30 - 9.30} \\
	\textbf{Key Results, Theorems, Definitions, etc.} \\
	\textbf{Charles Zitella}
\end{center}

\begin{abstract}
	This document contains a summary of all the key definitions, results, and theorems from class before the final exam. Point-set topology in Euclidean space; continuity and differentiability of functions in several variables. Implicit and inverse function theorems, Lagrange multipliers,
    Integrals in several variables, Fubini, change of variables, line and surface integrals, Green's theorem, etc...
\end{abstract}

\tableofcontents

\section{Continuity and Differentiability}

\subsection{Differentiability}

\begin{definition}[Differentiation]
    $f : \Omega \to \R^m$ where $\Omega \subset \R^n$ is a domain\footnote{Open, Connected}, 
    $f := (f_1, \dots, f_m), 
    f_i : \Omega\subset \R^n \to \R$ is \textbf{differentiable} at $x_0 \in \Omega$ if 
    there exists a linear map $L : \R^n \to \R^m$ such that
    \[
        \lim_{x \to x_0} \frac{\|f(x) - f(x_0) - L(x - x_0)\|}{\|x - x_0\|} = 0
    \] 
    Equivalently, $\forall \epsilon > 0 \; \exists \delta > 0$ such that for any $x \in \R$
    with $\|x - x_0\| < \delta$ \footnote{Here, $\| \cdot \|$ is the standard Euclidean norm on
     $\R^m$ and $\R^n$.} we have
    \[
        \| f(x) - f(x_0) -L(x - x_0)\| \leq \epsilon \|x - x_0\|
    \]   
\end{definition}

\begin{theorem}
    If $f : \Omega \to \R^m$ is differentiable at $x_0 \in \Omega$, then the linear map $L$ is unique.
\end{theorem}

\begin{definition}
    If $f$ differentiable at $x_0 \in \R$, with $f : \Omega \subset \R^n \to \R^m$ we write the 
    differential map of $f$ at $x_0 \in \Omega$ as $L = Df(x_0)$      
\end{definition}

\begin{prop}
    If $f$ is differentiable at $x_0 \in \Omega$, then $f$ is continuous at $x_0 \in \Omega$.
    In fact, it is Lipschitz continuous at $x_0$, i.e. $\exists k > 0$ and $\delta > 0$
    such that if $\| x - x_0 \| < \delta$, then $\| f(x) - f(x_0) \| \leq k \|x - x_0\|$       
\end{prop}

\begin{prop}
    If $f : \Omega \subset \R^n \to \R^m$ differentiable at $x_0 \in \Omega$, and 
    $f = (f_1, \dots, f_m)$, $f_i : \Omega \to \R$, then $f$ differentiable at $x_0 \in \Omega$
    iff $f_j : \Omega \to \R$ is differentiable at $x_0 \in \Omega$ for all $j = 1, \dots, m$.        
\end{prop}

\begin{definition}[Partial Derivatives]
    Let $f : \Omega \subset \R^n \to \R^m$, $f(x_1, \dots, x_m) = (f_1(x_1, \dots, x_n), \dots, 
    f_m(x_1, \dots, x_n))$.
    Then
    \[
    \pdv{f_j}{x_i} (x_0) := \lim_{h \to 0} \frac{[f_j(x_1, \dots, x_i + h, \dots, x_n) - 
    f_j(x_1, \dots, x_i, \dots, x_n)]}{h}
    \]
    if the limit exists. We call $\pdv{f_j}{x_i} (x_0)$ the \textbf{partial derivative} of $f_j$ 
    with respect to $x_i$ at $x_0$.
\end{definition}

\begin{prop}
    Let $f: \Omega \subset \R^m \to \R^m$, $f = (f_1, \dots, f_m)$, $f_j : \Omega \to \R$, 
    be differentiable at $x_0 \in \Omega$. Then all partial derivatives $\pdv{f_j}{x_i} (x_0)$ exist and
    \[
        L = Df(x_0) = \begin{bmatrix}
            \pdv{f_1}{x_1} (x_0) & \pdv{f_1}{x_2} (x_0) & \dots & \pdv{f_1}{x_n} (x_0) \\
            \pdv{f_2}{x_1} (x_0) & \pdv{f_2}{x_2} (x_0) & \dots & \pdv{f_2}{x_n} (x_0) \\
            \vdots & \vdots & \ddots & \vdots \\
            \pdv{f_m}{x_1} (x_0) & \pdv{f_m}{x_2} (x_0) & \dots & \pdv{f_m}{x_n} (x_0)
        \end{bmatrix}
    \]   
    called the Jacobian of $f$ at $x_0$\footnote{also called the derivative of $f$ at $x_0\in \Omega$}.  
\end{prop}

\begin{remark}
    $f$ differentiable at $x_0 \in \Omega \Rightarrow \pdv{f_j}{x_i} (x_0)$ exists for all $i, j$.
    However, the converse need not be true, for example the classic example of $f : \R^2 \to \R$,
    \[
        f(x, y) = \begin{cases}
            \frac{x^2 y}{x^2 + y^2} & (x, y) \neq (0, 0) \\
            0 & (x, y) = (0, 0)
        \end{cases}
    \]    
    Both partial derivatives of $f$ at $(0, 0)$ exist and are equal to $0$, but $f$ is not 
    differentiable at $(0, 0)$.

    Steps to solve: Find candidate for $Df(0, 0)$ using the partial derivatives, then show
    that the limit in the definition of differentiability does not go to $0$ for this candidate. 
\end{remark}

\begin{theorem}
    Let $f = (f_1, \dots, f_m) : \Omega \subset \R^n \to \R^m$. Suppose that all partial
    derivatives $\pdv{f_j}{x_i}$ exist and are continuous at $x_0 \in \Omega$. 
    Then $f$ is differentiable at $x_0$.
\end{theorem}

\begin{definition}
    Suppose $f : \Omega \subset \R^n \to \R$ has continuous partial derivatives 
    $\pdv{f}{x_i}$, $i = 1, \dots, n$ at all points $x \in \Omega$. Then $f$ is 
    continuously differentiable on $\Omega$, and we write $f \in C^1(\Omega)$.     
\end{definition}

\begin{remark}
    Continuity of partial derivatives is a sufficient condition for differentiability, but 
    it is not necessary. For example, the function $f : \R^2 \to \R$
    \[
        f(x, y) = \begin{cases}
            \frac{x^2y^2}{x^2 + y^2} & (x, y) \neq (0, 0) \\
            0 & (x, y) = (0, 0)
        \end{cases}
    \]  
    Then $\pdv{f}{x}(0, 0) = \pdv{f}{y}(0, 0) = 0$, however, along the line $x = y^2, x = t, y = t^2$,
    $\pdv{f}{x}(t^2, t) = \frac{1}{2} \neq 0$. But, $f$ is differentiable at 
    $(0, 0)$ with $Df(0, 0) = 0$.
\end{remark}

\begin{prop}[Basic Properties of Differentiation]
    \leavevmode
    \begin{enumerate}
        \item $f, g : \Omega \subset \R^n \to \R^m$ both differentiable at $x_0 \in \Omega$,
        then so is $F = f + g$. Moreover, $DF(x_0) = Df(x_0) + Dg(x_0)$.
        \item $f, g : \Omega \subset \R^n \to \R$ both differentiable at $x_0 \in \Omega$,
        then so is $F = f \cdot g : \Omega \to \R$. Moreover,
        $DF(x_0) = f(x_0) \cdot Dg(x_0) + g(x_0) \cdot Df(x_0)$.
        \item  $f, g : \Omega \subset \R^n \to \R$ both differentiable at $x_0 \in \Omega$,
        and $g(x_0) \neq 0$, then so is $F = \frac{f}{g} : \Omega \to \R$. Moreover,
        \[
            DF(x_0) = \frac{g(x_0) \cdot Df(x_0) - f(x_0) \cdot Dg(x_0)}{g(x_0)^2}
        \]  
        \item $f : \Omega \subset \R^n \to \tilde{\Omega} \subset \R^m$, 
        $g : \tilde{\Omega} \subset \R^m \to \R^k$ with $f$ differentiable at $x_0 \in \Omega$,
        and $g$ differentiable at $y_0 = f(x_0) \in \tilde{\Omega}$. Then 
        $H = g \circ f$ is differentiable at $x_0 \in \Omega$ and
        \[
            DH(x_0) = Dg(y_0) \cdot Df(x_0)
        \]       
    \end{enumerate}
\end{prop}

\begin{theorem}[Chain Rule]
    Let $f : \Omega \subset \R^n \to \tilde{\Omega} \subset \R^m$, $\tilde{\Omega} = f(\Omega)$,
    $g : \tilde{\Omega} \subset \R^m \to R^k$. If $f$ is differentiable at $x_0 \in \Omega$ and 
    $g$ is differentiable at $y_0 = f(x_0) \in \tilde{\Omega}$, then 
    $H = g \circ f : \Omega \subset \R^n \to R^k$ is differentiable at $x_0 \in \Omega$, and
    \[
        DH(x_0) = Dg(y_0) \cdot Df(x_0)
    \]      
\end{theorem}

\begin{exmp}
    Given $f$ differentiable in $R^2$ and $g(r, \theta) = (r\cos \theta, r\sin \theta)$, 
    with $(r, \theta) \in (0, \infty) \times [0, 2\pi]$. Set $F(r, \theta) = f(g(r, \theta))$,
    compute $\pdv{F}{r}(r, \theta)$ and $\pdv{F}{\theta}(r, \theta)$ in terms of the partial 
    derivatives of $f$, using the chain rule.

    To solve, the multivariable chain rule gives us
    \[
        \pdv{F}{r} = \pdv{f}{x} \cdot \pdv{x}{r} + \pdv{f}{y} \cdot \pdv{y}{r} = \pdv{f}{x} \cos \theta + \pdv{f}{y} \sin \theta
    \] 
    and
    \[
        \pdv{F}{\theta} = \pdv{f}{x} \cdot \pdv{x}{\theta} + \pdv{f}{y} \cdot \pdv{y}{\theta} = -r \pdv{f}{x} \sin \theta + r \pdv{f}{y} \cos \theta
    \] 

\end{exmp}

\subsection{Tangent Planes and Directional Derivatives}

\begin{definition}[Gradient]
Let $f : \Omega \subset \R^n \to \R$ be differentiable on $\Omega$. Then 
$Df(x) = (\pdv{f}{x_1}(x), \dots, \pdv{f}{x_n}(x)) = \nabla f(x)$ is called the \textbf{gradient} of 
$f$ at $x \in \Omega$. 
\end{definition}

\begin{definition}[Tangent Planes]
Let $f : \Omega \subset \R^n \to \R$ be differentiable on $\Omega$. 
$S := \operatorname{graph}f = \left\{ (x, z) \in \Omega \times \R : z = f(x) \right\} $.
Then for $x_0 \in \Omega$, $T_{x_0}S = \left\{ (x, z) \in \R^n \times \R :
z = f(x_0) + Df(x_0) \cdot (x - x_0) \right\} $  
\end{definition}

\begin{definition}[Directional Derivative]
    Let $f : \Omega \subset \R^n \to \R$ be differentiable on $\Omega$. Then for $x_0 \in \Omega$ and 
    $v \in \R^n$, the directional derivative of $f$ at $x_0$ in the direction of $v$ is defined as
    \[
        D_vf(x_0) = \langle \nabla f(x_0), v \rangle
    \] 
    
\end{definition}

\begin{prop}
    Suppose $\nabla f(x) \neq 0$. Then $\nabla f(x)$ points in the direction of the steepest
    increase of $f$. \footnote{Proof follows from the Cauchy-Schwarz inequality, $<a, b> = \|a\| \cdot \|b\| \cos \theta$  }   
\end{prop}

\subsection{Clairaut's Theorem}

\begin{theorem}[MVT in $\R^n$]
    Let $B \subset \R^n$ be a ball, $B = \left\{ x \in \R^n  : \| x - x_0 \| < p \text{ for some } x_0
    \right\} $, and $f : B \to \R$ differentiable on $B$. Then, for any $x, y \in B$ there exists
    $z \in B$ such that     
    \[
        f(y) - f(x) = Df(z) \cdot (y - x)
    \] 
    As a consequence, $\| f(y) - f(x) \| \leq \| Df(z) \| \cdot \| y - x \|$.   
\end{theorem}

\begin{definition}
We say $f$ twice differentiable at $x \in \Omega$ if $Df$ exists locally near $x \in \Omega$ and
$Df : \Omega \subset \R^n \to \R^{n \times m}$ is differentiable at $x \in \Omega$. Then
$D^2f = D(Df)$ at $x \in \Omega$. Similarly, for all $k \geq 2$, $D^kf = D(D^{k - 1}f)$.          
\end{definition}

\begin{definition}
    Given $f : \Omega \subset \R^n \to \R$, we say that $f \in C^k(\Omega)$ if all the partial
    derivatives up to order $k$ are continuous on $\Omega$.    
\end{definition}

\begin{definition}(Hessian)
    Given $f : \Omega \subset \R^n \to \R$, the \textbf{Hessian} matrix of $f$ at $x \in \Omega$ is
    \[
        Hf(x) = \begin{bmatrix}
            \pdv{^2f}{x_1 \partial x_1} (x) & \pdv{^2f}{x_1 \partial x_2} (x) & \dots & \pdv{^2f}{x_1 \partial x_n} (x) \\
            \pdv{^2f}{x_2 \partial x_1} (x) & \pdv{^2f}{x_2 \partial x_2} (x) & \dots & \pdv{^2f}{x_2 \partial x_n} (x) \\
            \vdots & \vdots & \ddots & \vdots \\
            \pdv{^2f}{x_n \partial x_1} (x) & \pdv{^2f}{x_n \partial x_2} (x) & \dots & \pdv{^2f}{x_n \partial x_n} (x)
        \end{bmatrix}
    \]    

    For $f(x, y)$, $Df = (\pdv{f}{x}, \pdv{f}{y})$ and 
    \[
        D^2f = Hf = \begin{bmatrix}
            \pdv{^2f}{x \partial x} & \pdv{^2f}{x \partial y} \\
            \pdv{^2f}{y \partial x} & \pdv{^2f}{y \partial y}
        \end{bmatrix}
    \]    
\end{definition}

\begin{theorem}[Clairaut's Theorem]
    Let $f : \Omega \subset \R^n \to \R$ be twice differentiable at $x \in \Omega$. Then,
    \[
        \pdv{^2f}{x_i \partial x_j} (x) = \pdv{^2f}{x_j \partial x_i} (x)
    \]   
    for all $i, j = 1, \dots, n$. Moreover, if $f \in C^2(\Omega)$, then the above holds for all
    $x \in \Omega$. 
\end{theorem}

\begin{corollary}
If $\pdv{^2f}{x_i \partial x_j} (x)$ are all continuous at $x \in \Omega$ for all $i, j = 1, \dots, n$, 
then $\pdv{^2f}{x_i \partial x_j} (x) = \pdv{^2f}{x_j \partial x_i} (x)$ for all $i, j = 1, \dots, n$.  
\end{corollary}

\subsection{Inverse Function Theorem and Implicit Function Theorem}

\begin{theorem}[Banach Fixed Point]
Let $(X, d)$ be a complete metric space, and $f : X \to X$ be a contraction mapping with
$d(f(x), f(y)) \leq \alpha d(x, y)$ for some $0 < \alpha < 1$. Then $\exists$ unique 
$x_0 \in X$ such that $f(x_0) = x_0$.       
\end{theorem}

\begin{definition}
    Let $M_n := \left\{ n \times n \text{ matrices} \right\} \cong \R^{n^2}$ and 
    $\| A \| := \sqrt{\sum_{i,j} a_{ij}^2}$ for $A = (a_{ij}) \in M_n$.

    Let $GL_n := \left\{ A \in M_n : \det A \neq 0 \right\}  = \det ^{-1}(\R \setminus \{0\})$ 
    Thus since $\det : M_n \to \R$ is continuous, $GL_n$ is an open subset of $M_n$.

\end{definition}

\begin{theorem}[Inverse Function]
    Let $f : \Omega \subset \R^n \to \R^n$ be $C^1$. Let $x_0 \in \Omega$ and assume $Df(x_0) 
    \in GL_n$. Then there exists neighbourhoods $U$ and $V$ of $x_0$ and $f(x_0)$ respectively
    such that $f(U) = V$ and $f|_{U}$ has a $C^1$ inverse map $f^{-1} : V \to U$.
    Moreover, for any $y \in V$, and $x = f^{-1}(y)$, we have $Df^{-1}(y) = [Df(x)]^{-1}$.             
\end{theorem}

\begin{remark}
    If $f \in C^k$, $k \geq 1$, then the inverse map $f^{-1}$ is also in $C^k$.  
\end{remark}

\begin{exmp}
    IFT is local. For example, $f(x, y) = (e^y \cos x, e^y \sin x)$. Then 
    $Df(x, y) = \begin{bmatrix}
        -e^y \sin x & e^y \cos x \\
        e^y \cos x & e^y \sin x
    \end{bmatrix}$, and $\det Df(x, y) = -(\cos ^2 x + \sin ^2 x) e^{2y} = -e^{2y} \neq 0$ for all 
    $(x, y) \in \R^2$. However, $f(x + 2k \pi, y) = f(x, y)$ for all $k \in \mathbb{Z}$, so $f$ is not 
    injective globally on $(x, y)$. 
\end{exmp}

\begin{theorem}[Implicit Function Theorem]
    $F : \Omega \subset \R^n \times \R^m \to \R^m$ $C^k$ map, $X = (x, y) \in R^n \times \R^m$ and
    $X_0 = (x_0, y_0)$ with $F(X_0) = 0$. Writing $F = (F_1, \dots, F_m)$, assume that
    \[
        D_yF(X_0) = \begin{bmatrix}
            \pdv{F_1}{y_1} & \pdv{F_1}{y_2} & \dots & \pdv{F_1}{y_m} \\
            \pdv{F_2}{y_1} & \pdv{F_2}{y_2} & \dots & \pdv{F_2}{y_m} \\
            \vdots & \vdots & \ddots & \vdots \\
            \pdv{F_m}{y_1} & \pdv{F_m}{y_2} & \dots & \pdv{F_m}{y_m} 
        
        \end{bmatrix} (X_0) \in GL_m
    \]       
    Then there exists neighbourhoods $U$ and $V$ of $x_0 \in \R^n$ and $y_0 \in \R^m$ and 
    a unique $C^k$ map $f : U \to V$ such that $F(x, f(x)) = 0$, for all $x \in U$.         
\end{theorem}

\begin{corollary}
    Let $F : \Omega \subset \R^n \to \R$ be $C^k(\Omega)$. Let $x = (x', y) \in \R^{n - 1} \times \R$
    and suppose $(x_0', y_0) \in \Omega$ with $\pdv{f}{y}(x_0', y_0) \neq 0$. Then there exists
    neighbourhoods $U$, $V$ of $x_0' \in \R^{n -1}$ and $y_0 \in \R$ and a unique function 
    $f : U \to V$, $f \in C^k(U)$ such that $\left\{ F(x', y) = 0 \right\} = \left\{ y = f(x') \right\}
    \forall x' \in U$.            
\end{corollary}

\begin{theorem}[Morse Lemma]
    $f : \Omega \subset \R^n \to \R$ a $C^k$ function $k \geq 3$. Let $0 \in \Omega$ be a 
    critical point, i.e. $\nabla f(0) = 0$, and assume $f(0) = 0, \nabla f(0) = 0, Hf(0) \in GL_n$. 
    Then there exists open sets $U$ and $V$ of $0 \in U \cap V$ and $g \in C^{k -2}(U)$, $g : U \to V$
    with $g^{-1} : U \to V$, $g^{-1} \in C^2(V)$ such that $f(g(y)) = 
    y_{\ell + 1}^2 + \dots + y_n^2 - (y_1^2 + y_2^2 + \dots y_\ell^2)$. Here $\ell \in \mathbb{Z}$,
    $\ell \in [0, n]$.            
\end{theorem}

\subsection{Taylor's Theorem}

\begin{theorem}[Taylor's Theorem in $\R^n$ ]
    Let $f : \Omega \subset \R^n \to \R$, $f \in C^{k + 1}(\Omega)$. Then, for any 
    $x^0 \in \Omega$,
    \[
        f(x) = f(x^0) + \sum_{j = 1}^k \frac{1}{j!} \left[\sum_{i_1 = 1}^n \dots \sum_{i_j = 1}^n 
        \pdv{^jf}{x_{i_1} \dots \partial x_{i_j}} (x^0) (x_{i_1} - x_{i_1}^0) 
        \dots (x_{i_j} - x_{i_j}^0)\right] + R_k(f)(x)
    \]  
    where $\frac{|R_k(f)(x)|}{\| x - x^0 \|^k } \leq M_k \| x - x^0 \| $ as $x \to x^0$.  
\end{theorem}

\begin{remark}
    The k-th order Taylor polynomial of $f$ at $x^0$ is the sum of the first two terms 
    in the above expansion, i.e. \[P_k(f)(x) = f(x^0) + \sum_{j = 1}^k \frac{1}{j!} 
    \left[\sum_{i_1 = 1}^n \dots \sum_{i_j = 1}^n 
        \pdv{^jf}{x_{i_1} \dots \partial x_{i_j}} (x^0) (x_{i_1} - x_{i_1}^0) 
        \dots (x_{i_j} - x_{i_j}^0)\right]\]
         This is a good approximation to $f(x)$ near $x^0$
        provided $\lim_{x \to x^0} \frac{|R_k(f)(x)|}{\| x - x^0 \|^k } = 0$.   
\end{remark}

\begin{theorem}
    Let $f : \Omega \subset \R^n \to \R$ be $C^{k-1}(\Omega)$ and assume $D^{k-1}f$ is differentiable
    at $x^0 \in \Omega$. Then the $k$-th order Taylor expansion in the previous theorem still holds,
    but with the only assumption of the remainder term being
    $\lim_{x \to x^0} \frac{|R_k(f)(x)|}{\| x - x^0 \|^k } = 0$. 
\end{theorem}

\begin{exmp}
    Let $f(x, y) = e^x + sin(xy)$. Compute the second order Taylor expansion of $f$ around $(0, 0)$.
    
    To solve this problem, we first recognize that $f \in C^\infty(\R^2)$, 
    so we can apply Taylor's theorem. We compute the first and second order partial derivatives of 
    $f$ at $(0, 0)$:
    \[
        f_x = \pdv{f}{x} = e^x + y\cos(xy), \quad f_y = \pdv{f}{y} = x\cos(xy)
    \] 
    Therefore $f_x(0, 0) = 1$ and $f_y(0, 0) = 0$. Next, we compute the second order partial derivatives:
    \[
        f_{xx} = \pdv{^2f}{x^2} = e^x - y^2\sin(xy), \quad f_{yy} = \pdv{^2f}{y^2} = -x^2\sin(xy), 
        \quad f_{xy} = \pdv{^2f}{x\partial y} = \pdv{^2f}{y\partial x} = \cos(xy) - xy\sin(xy)
    \] 
    Therefore, $f_{xx}(0, 0) = 1$, $f_{yy}(0, 0) = 0$, and $f_{xy}(0, 0) = 1$. 
    Thus, the second order Taylor expansion of $f$ around $(0, 0)$ is
    \[
        f(x, y) = f(0, 0) + f_x(0, 0)x + f_y(0, 0)y + \frac{1}{2!}\left[f_{xx}(0, 0)x^2 + 
        2f_{xy}(0, 0)xy + f_{yy}(0, 0)y^2\right] + R_2(f)(x, y)
    \] 
    which simplifies to
    \[
        f(x, y) = 1 + x + \frac{x^2}{2} + xy + R_2(f)(x, y)
    \] 
    and the remainder term satisfies $\frac{|R_2(f)(x, y)|}{x^2 + y^2} \leq M \sqrt{x^2 + y^2}$ as 
    $(x, y) \to (0, 0)$.
\end{exmp}

\subsection{Lagrange Multipliers}

\begin{theorem}[Lagrange Multiplier]
    Let $f, g : \Omega \subset \R^n \to \R$ be $C^1$ functions in $\R$. Let $\Sigma = 
    \left\{ x \in \Omega : g(x) = 0 \right\} $ and assume that $Dg(x_0) \neq 0$ for all
    $x_0 \in \Sigma$. Then if $f|_\Sigma$ has a max or min at $x_0 \in \Sigma$, then there
    exists $\lambda \in \R$ such that 
    \[
        \nabla f(x_0) = \lambda g(x_0)
    \]         
\end{theorem}

\begin{theorem}[Max/Min]
    Let $\tilde{\Omega} \supset \Omega$ be a slightly larger domain and assume $f, g \in C^1(\Omega)$,
    where $\partial \Omega = \left\{ x \in \tilde{\Omega} : g(x) = 0 \right\} $, $\nabla g(x) \neq 0
    \forall x \in \partial \Omega$.
    Then,
    \begin{enumerate}
        \item If $x_0 \in \partial \Omega$ is a local max/min of $f|_{\partial \Omega}$, then
        $\nabla f(x_0) = \lambda \nabla g(x_0)$ for some $\lambda \in \R$.
        \item If $x_0 \in \operatorname{Int}(\Omega)$ is a local max/min of $f$ in 
        $\operatorname{Int}(\Omega)$, then $\nabla f(x_0) = 0$.
        \item Global max/min of $f|_{\tilde{\Omega}}$ is largest/smallest value of $f$ at all points
        in 1 and 2.         
    \end{enumerate}   
\end{theorem}

\begin{theorem}[Multiple Constraints]
    $f : \Omega \subset \R^n \to \R$, $g_i : \Omega \subset \R^n \to \R$, $i = 1, \dots, k$ are $C^1$ 
    functions. Let $\Gamma_0 = \left\{ x \in \Omega : g_i(x) = 0, i = 1, \dots, k \right\} $
    and assume that $\left\{ \nabla g_1(x), \dots, \nabla g_k(x) \right\} $ are linearly
    independent for all $x \in \Gamma_0$. Then if $f_{\Gamma_0}$ has a max/min at
    $x_0 \in \Gamma_0$, then there exists $\lambda_1, \dots, \lambda_k \in \R$ such that
    \[
        \nabla f(x_0) = \sum_{i = 1}^k \lambda_i \nabla g_i(x_0)
    \]     
    
\end{theorem}

\begin{exmp}
    Find the absolute max and min of $f(x_1, \dots, x_n) = x_1 + x_2$ on 
    $B_1(0) = \left\{ |x|^2 \leq 1 \right\} = \left\{ x_1^2 + \dots + x_n^2 \leq 1 \right\} $,
    $\partial B_1(0) = \left\{ x_1^2 + \dots + x_n^2 = 1 \right\} $.
    
    To solve, first observe $f \in C^k(\R^n)$ for all $k \geq 1$. Thus we will apply the max/min
    theorem with $g(x) = |x|^2 - 1 = x_1^2 + \dots + x_n^2 - 1$. Then, since
    $\nabla f(x) = (1, 1, 0, \dots, 0) \neq 0$, there are no critical points in $\operatorname{Int}
    (B_1(0)) = \left\{ |x|^2 < 1  \right\} $. So the global max/min must occur on 
    $\partial B_1(0)$.
    
    By Lagrange multiplier theorem, at max/min $x_0 \in \partial B_1(0)$, 
    there exists $\lambda \in \R$ such that 
    \begin{align}
        \nabla f(x_0) = \lambda \nabla g(x_0) &\Leftrightarrow (1, 1, 0, \dots, 0) = \lambda 
        (2x_{1}^0, 2x_{2}^0, 2x_{3}^0, \dots, 2x_{n}^0) \\
        &\Leftrightarrow x_{1}^0 = x_{2}^0 = \frac{1}{2\lambda}, x_{3}^0 = \dots = x_{n}^0 = 0
    \end{align}
    Also, 
    \[
        (x_1^0)^2 + (x_2^0)^2 = 1 \Leftrightarrow \frac{1}{4\lambda^2} + \frac{1}{4\lambda^2} = 1
        \Leftrightarrow \lambda = \pm \frac{1}{\sqrt{2}}
    \] 

    Thus, the candidates for max/min are $x_0 = \left( \frac{1}{\sqrt{2}}, \frac{1}{\sqrt{2}}, 0, \dots, 0 \right) $
    and $x_0 = \left( -\frac{1}{\sqrt{2}}, -\frac{1}{\sqrt{2}}, 0, \dots, 0 \right) $. Evaluating $f$ at these points, we get
    \[
        f\left( \frac{1}{\sqrt{2}}, \frac{1}{\sqrt{2}}, 0, \dots, 0 \right) = \sqrt{2}, \quad f\left( -\frac{1}{\sqrt{2}}, -\frac{1}{\sqrt{2}}, 0, \dots, 0 \right) = -\sqrt{2}
    \]
    Thus, the absolute max of $f$ on $B_1(0)$ is $\sqrt{2}$ and the absolute min is $-\sqrt{2}$.
\end{exmp}


\section{Riemann Integration}
\subsection{Integration over Box domains}

\begin{definition}[Partition]
    Let $B = [a_1, b_1] \times \dots \times [a_n, b_n]$. A \textbf{partition} of $B$ is an ordered set 
    \[
        \mathcal{P} = \left\{ P_1 = \{ x_1^0, x_1^1, \dots, x_1^n\}, \dots, P_n = \{ x_n^0, x_n^1, \dots, x_n^m \} \right\} \quad \footnote{The lower index denotes the dimension, and the upper index denotes the parition index.} 
    \]   
    Here, for any given $j$, $P_j$ is the set of $x_j^i$ parition of $[a_j, b_j]$  s.t.  
    \[
        a_j = x_j^0 < x_j^1 < \dots < x_j^{m_j - 1} < x_j^{m_j} = b_j
    \]  
\end{definition}
\begin{definition}[Refinement]
    Given two partitions $\mathcal{P}$ and $\tilde{\mathcal{P}}$, we say that $\tilde{\mathcal{P}}$ is a \textbf{refinement} of $\mathcal{P}$,
    i.e. $\tilde{\mathcal{P}} \supset \mathcal{P}$, if for each $j = 1, \dots, n$,
    \[
        P_j = \left\{ x_j^0, x_j^1, \dots, x_j^{m_j} \right\} \subset \left\{ \tilde{x}_j^0, \tilde{x}_j^1, \dots, \tilde{x}_j^{m_j} \right\} = \tilde{P}_j 
    \]       
\end{definition}
\begin{definition}[Sub-box]
    Given partition $\mathcal{P}$ of $B$, $\mathcal{P} = \left\{ P_1, \dots, P_j \right\} $, let
    \[
        B_{i_1, i_2, \dots, i_n} = [x_1^{i_1 - 1}, x_1^{i_1}] \times \dots \times [x_n^{i_n - 1}, x_n^{i_n}] 
    \]    
    be a \textbf{sub-box} of $B$ for some $\{i_1, \dots, i_n\}$ partition indices, i.e. $1 \leq i_j \leq m_j$ for all $j = 1, \dots, n$.
    Then define the length of the $j$th side of the sub-box as 
    \[
        \Delta x_j^{i_j} := x_j^{i_j} - x_j^{i_j - 1} \quad \forall j
    \] 
    Thus the volume is intuitively defined as
    \[
    \Delta V_{i_1, \dots, i_n} := \Delta x_1^{i_1} \cdot \dots \cdot \Delta x_n^{i_n} = \prod_{j = 1}^n (x_j^{i_j} - x_j^{i_j - 1})
    \] 

\end{definition}

\begin{definition}[Upper and Lower Riemann Sums]
    Let $f : B \to \R$ be bounded and $\mathcal{P}$ a partition of $B$. Define the \textbf{lower} and \textbf{upper} Riemann sums of $f$ with respect to $\mathcal{P}$ as
    \[
        L(f, \mathcal{P}) = \sum_{i_1 = 1}^{m_1} \dots \sum_{i_n = 1}^{m_n} \left(\inf_{x \in B_{i_1, \dots, i_n}} f(x)\right) \cdot \Delta V_{i_1, \dots, i_n}
        \quad \text{and} \quad U(f, \mathcal{P}) = \sum_{i_1 = 1}^{m_1} \dots \sum_{i_n = 1}^{m_n} \left(\sup_{x \in B_{i_1, \dots, i_n}} f(x)\right) \cdot \Delta V_{i_1, \dots, i_n}
    \]    
\end{definition}

\begin{lemma}
    Let $f : B \to \R$ be bounded, and $\tilde{\mathcal{P}} \supset \mathcal{P}$ be a refinement. Then
    \[
        L(f, \mathcal{P}) \leq L(f, \mathcal{\tilde{P}}) \leq U(f, \mathcal{\tilde{P}}) \leq U(f, \mathcal{P})
    \] 
\end{lemma}

\begin{definition}
    Let $f : B \to \R$ be bounded, i.e. $\exists m, M$ s.t. $m \leq f(x) \leq M$, $\forall x \in \R$. Then define
    \[
        U(f) := \inf_{\mathcal{P}} U(f, \mathcal{P}) \quad \text{and} \quad L(f) := \sup_{\mathcal{P}} L(f, \mathcal{P})
    \]   
\end{definition}

\begin{definition}
    Let $f : B \to \R$ be bounded. Then we say $f$ \textbf{Riemann integrable} over $B$ provided $L(f) = U(f)$, and we define
    \[
        \int_B f(x) \; dV := L(f) = U(f)
    \]   
\end{definition}

\begin{lemma}
    Let $f : B \to \R$ be bounded. Then $f$ is Riemann integrable over $B$ iff $\forall \epsilon > 0$, $\exists \mathcal{P}$ partition
    of $B$ such that $U(f, \mathcal{P}) - L(f, \mathcal{P}) \leq \epsilon$.    
\end{lemma}

\begin{theorem}
    Let $f : B \to \R$ be $C(B)$. Then $f$ is Riemann integrable over $B$. (Continuity implies integrability)  
\end{theorem}

\begin{lemma}[MVT for Integrals]
    Let $f : B \to \R$ be continuous. Then $\exists x_0 \in B$ s.t. $\int_B f(x) \; dV = f(x_0) \cdot \Vol(B)$, where
    $\Vol(B) = \int_B 1 \; dV$. 
\end{lemma}

\subsection{Integration over General Domains}

\begin{definition}
    Given $\Omega \subset \R^n$ bounded, $f : \Omega \to \R$ bounded, we let $B \subset \R^n$ be a box.
    We define $F : B \to \R$ as 
    \[
        F(x) := \begin{cases}
            f(x) & x \in \Omega \\
            0 & x \in B \setminus \Omega
        \end{cases}
    \]     
    Then
    \[
        \int_{\Omega} f \; dV = \int_B F \; dV \quad \text{ when the latter exists}
    \] 
\end{definition}

\begin{definition}[Content Zero]
    $\Gamma \subset \R^n$ is of \textbf{Content 0} if $\forall \epsilon > 0$, there exists open box domains $B_1^0, \dots, B_N^0$
    s.t.
    \begin{enumerate}
        \item $\Gamma \subset \bigcup_{i = 1}^N B_i^0$
        \item $\sum_{i = 1}^N \Vol(B_i^0) < \epsilon$
    \end{enumerate}  
\end{definition}

\begin{exmp}
    If $\Gamma_1, \dots, \Gamma_n$ are all content 0, then $\Content(\Gamma_1 \cup \dots \cup \Gamma_N) = 0$  

    To solve, let $\epsilon > 0$. Since $\Gamma_i$ is of content 0 for all $i = 1, \dots, N$, there exists open box domains $B_1^i, \dots, B_{M_i}^i$ s.t. $\Gamma_i \subset \bigcup_{j = 1}^{M_i} B_j^i$ and $\sum_{j = 1}^{M_i} \Vol(B_j^i) < \frac{\epsilon}{N}$ for all $i$. Then we have
    \[
    \Gamma_1 \cup \dots \cup \Gamma_N \subset \bigcup_{i = 1}^N \bigcup_{j = 1}^{M_i} B_j^i
    \]
    and
    \[
    \sum_{i = 1}^N \sum_{j = 1}^{M_i} \Vol(B_j^i) < \sum_{i = 1}^N \frac{\epsilon}{N} = \epsilon
    \] 
\end{exmp}

\begin{theorem}
    Suppose $f : B \subset \R^n \to \R$ bounded. Let $\Gamma(f) = \text{discontinuity set of } f$. 
    
    Then $\Content(\Gamma(f)) = 0$ iff $f$ is Riemann integrable over $B$.   
\end{theorem}

\begin{exmp}
    $\Omega \subset \R^{n + 1}$ compact. $g : \Omega \to \R $ continuous. Then the graph of $g$, 
    $\Gamma = \left\{ (x', g(x')) : x' \in \Omega \right\} $, satisfies $\Content(\Gamma) = 0$.
    
    To solve, 
\end{exmp}

\begin{corollary}
    Suppose $f : B \to \R$ is bounded and $\Gamma(f)$ is the union of finitely many graphs of continuous functions over closed
    domains in $\R^{n-1}$. Then $f$ is Riemann integrable over $B$.   
\end{corollary}

\begin{corollary}
    $\Omega \subset \R^n$ closed and bounded, $\partial \Omega$ is the union of finitely many graphs of
    $C^0$ functions over closed domains in $\R^{n - 1}$. Suppose $f : \Omega \to \R$ bounded and
    $\Content(\Gamma(f) \cap \Omega) = 0$. Then $f$ is Riemann integrable over $\Omega$.      
\end{corollary}

\begin{corollary}
    Suppose $\Omega$ closed and bounded, $\Omega \subset \tilde{\Omega}$ with $G \in C^1(\tilde{\Omega})$ s.t.
    $\Omega = \{G(x) \leq 0 : x \in \tilde{\Omega} \}$ and $\partial \Omega = \{G(x) = 0 \}$ with $\nabla G(x) \neq 0$ 
    for all $x \in \partial \Omega$. Suppose $f : \Omega \to \R$ is bounded and $\Content(\Gamma(f) \cap \Omega) = 0$.
    Then $f$ is Riemann integrable over $\Omega$.  
\end{corollary}


\subsection{Basic Properties of Integrals}

\begin{prop}
    Assume $\Omega \subset \R^n$ bounded and $f, g : \Omega \to \R$ are both Riemann integrable over $\Omega$.
    Then,
    \begin{enumerate}
        \item $\int_{\Omega} (f + g) \; dV = \int_{\Omega} f \; dV + \int_{\Omega} g \; dV$;
        \item $\int_{\Omega} \alpha f \; dV = \alpha \int_{\Omega} f \; dV$ for any $\alpha \in \R$;
        \item $f \leq g \quad \Rightarrow \quad \int_{\Omega} f \; dV \leq \int_{\Omega} g \; dV$;
        \item If $\Omega = \Omega_1 \cup \Omega_2$, $\Omega$ bounded and $\Content(\Omega_1 \cap \Omega_2) = 0$.
        Then $\int_{\Omega} f \; dV = \int_{\Omega_1} f \; dV + \int_{\Omega_2} f \; dV$.     
    \end{enumerate}   
\end{prop}

\subsection{Fubini's Theorem}

\begin{definition}
    Given $f : B \to \R$ bounded, we define the iterated integrals as
    \[
        \int_a^b \int_c^d f(x,y) \; dy \; dx = \int_a^b \left( \int_c^d f(x,y) \; dy \right) dx \quad \text{ and } \quad \int_c^d \int_a^b f(x,y) \; dx \; dy = \int_c^d \left( \int_a^b f(x,y) \; dx \right) dy
    \] 
\end{definition}

\begin{theorem}[Fubini for box domains]
    Suppose $f : B \to \R$ is Riemann integrable over $B$, where $B = [a,b] \times [c,d]$.
    \begin{enumerate}
        \item Assume $\int_c^d f(x,y) \; dy$ exists $\forall x \in [a,b]$. Then $\int_a^b \int_c^d f(x,y) \; dy \; dx$ exists and equals $\int_B f(x,y) \; dV$.
        \item Assume $\int_a^b f(x,y) \; dx$ exists $\forall y \in [c,d]$. Then $\int_c^d \int_a^b f(x,y) \; dx \; dy$ exists and equals $\int_B f(x,y) \; dV$.  
    \end{enumerate}   
    
\end{theorem}

\subsection{Fubini for General Domains}

\begin{definition}[y-simple, x-simple domains]
    A domain $\Omega \subset \R^2$ is called \textbf{$y$-simple} if it is of the form
    \[
        \Omega = \left\{ a \leq x \leq b : \varphi_1(x) \leq y \leq \varphi_2(x) \right\} 
    \]   
    where $\varphi_1, \varphi_2 \in C^0([a,b])$.
    
    Similarly, $\Omega \subset \R^2$ is called \textbf{$x$-simple} if it is of the form
    \[
        \Omega = \left\{ c \leq y \leq d : \psi_1(y) \leq x \leq \psi_2(y) \right\} 
    \]    
    where $\psi_1, \psi_2 \in C^0([c,d])$.
    $\Omega \subset \R^2$ is called \textbf{simple} if it is both $y$-simple and $x$-simple. 
\end{definition}

\begin{theorem}[Fubini]
    Suppose $\Omega \subset \R^2$ bounded and $f$ integrable over $\Omega$. Then,
    \begin{enumerate}
        \item If $\Omega$ is $y$-simple and $\forall x \in [a,b]$, $\int_{\varphi_1(x)}^{\varphi_2(x)} f(x,y) \; dy$ exists, then
        \[
            \int_{\Omega} f(x,y) \; dV = \int_a^b \int_{\varphi_1(x)}^{\varphi_2(x)} f(x,y) \; dy \; dx
        \] 
        \item If $\Omega$ is $x$-simple and $\forall y \in [c,d]$, $\int_{\psi_1(y)}^{\psi_2(y)} f(x,y) \; dx$ exists, then
        \[
            \int_{\Omega} f(x,y) \; dV = \int_c^d \int_{\psi_1(y)}^{\psi_2(y)} f(x,y) \; dx \; dy
        \]     
    \end{enumerate} 
    If $\Omega$ is simple, then both hold. 
\end{theorem}

\begin{definition}
    $\Omega \subset \R^2$ is \textbf{elementary} if $\Omega$ is $x$-simple or $y$-simple.
\end{definition}

\begin{exmp}
    Compute $\iint_{\Omega} x \; dA$ where $\Omega = \{(x,y) \in \R^2 : 0 \leq x \leq \sqrt{\pi}, 0 \leq y \leq \sin x^2 \}$
    
    To solve, first note that $f(x,y) = x$ is $C^0$ on $\Omega$. Since $\Omega$ is bounded, $\iint_{\Omega} x \; dV$ exists. Also,
    $\Omega$ is $y$-simple with $\varphi_1(x) = 0, \varphi_2(x) = \sin x^2$, $\varphi_1, \varphi_2 \in C^0([0, \sqrt{\pi}])$.
    By Fubini,
    \[
        \iint_{\Omega} x \; dA = \int_0^{\sqrt{\pi}} \int_0^{\sin x^2} x \; dy \; dx = \int_0^{\sqrt{\pi}} x \sin x^2 \; dx = \frac{1}{2} \int_0^{\pi} \sin u \; du = \frac{1}{2} \left[ -\cos u \right]_0^{\pi} = 1
    \]     
\end{exmp}

\begin{exmp}
    Compute $I = \int_0^1 \int_y^1 e^{x^2} \; dx \; dy$.

    To solve, first note $\int_y^1 e^{x^2} \; dx$ doesn't exist in closed-form. But since $e^{x^2} \in C^0(\Omega)$, $\iint_{\Omega} e^{x^2} \; dA$ exists.
    Since $\Omega$ is simple, by Fubini,
    \[
        \iint_{\Omega} e^{x^2} \; dA = \int_0^1 \int_0^x e^{x^2} \; dy \; dx = \int_0^1 x e^{x^2} \; dx = \frac{1}{2} \int_0^1 e^u \; du = \frac{1}{2}(e - 1)
    \]   
\end{exmp}

\begin{definition}
    Some notation: Given a box domain $B \subset \R^n$, write
    \[
        \underline{\int_{B}} f \; dV := L(f) = \sup_{\mathcal{P}} L(f, \mathcal{P}) \quad \text{and} \quad \overline{\int_{B}} f \; dV := U(f) = \inf_{\mathcal{P}} U(f, \mathcal{P})
    \]  
\end{definition}

\begin{theorem}
    Let $B = B_1 \times B_2$ box domain where $B_1 \subset \R^k$ and $B_2 \subset \R^{n - k}$ also a box.
    Suppose $f : B \to \R$ is Riemann integrable. Set
    \[
        \ell(x) = \underline{\int_{B_2}} f(x,y) \; dV(y) \quad u(x) = \overline{\int_{B_2}} f(x,y) \; dV(y) \quad \forall x \in B_1 \subset \R^k
    \]      
    and 
    \[
        \tilde{\ell}(y) = \underline{\int_{B_1}} f(x,y) \; dV(x) \quad \tilde{u}(y) = \overline{\int_{B_1}} f(x,y) \; dV(x) \quad \forall y \in B_2 \subset \R^{n-k}
    \] 
    Then $\ell(x), u(x)$ are Riemann integrable over $B_1$. Also $\tilde{\ell}(y), \tilde{u}(y)$ are Riemann integrable over $B_2$.
    Moreover,
    \begin{align*}
        \int_B f(x,y) \; dV &= \int_{B_1} \ell(x) \; dV(x) = \int_{B_1} u(x) \; dV(x) \\
        \int_B f(x,y) \; dV &= \int_{B_2} \tilde{\ell}(y) \; dV(y) = \int_{B_2} \tilde{u}(y) \; dV(y)
    \end{align*}
\end{theorem}

\begin{definition}[$z$-simple]
    A bounded domain $W \subset \R^3$ is called \textbf{$z$-simple} if $\exists $ an elementary region $\Omega \subset \R^2$
    and $\eta_1, \eta_2 \in C^0(\Omega)$ s.t.
    \[
        W = \{ (x,y,z) \in \R^3 : (x,y) \in \Omega  \text{ and } \eta_1(x,y) \leq z \leq \eta_2(x,y) \}
    \]       
    It is possible to permute the coordinates and define $x$-simple and $y$-simple domains in $\R^3$ similarly.
\end{definition}

\begin{theorem}[Fubini in $\R^3$]
    Suppose $W \subset \R^3$ is bounded, $z$-simple, and $f : W \to \R$ os Riemann integrable over $W$. Assume
    $\forall (x,y) \in \Omega$, 
    \[
        \int_{\eta_1(x,y)}^{\eta_2(x,y)} f(x,y,z) \; dz \quad \text{exists}
    \]     
    Then it is Riemann integrable over $\Omega$, and
    \[
        \iiint f \; dV = \int_{\Omega} \left( \int_{\eta_1(x,y)}^{\eta_2(x,y)} f(x,y,z) \; dz \right) dV(x,y)
    \]   
    The same holds for $x$-simple and $y$-simple.  
\end{theorem}

\subsection{Change of Variables}

\begin{definition}[Regular Partition]
    Let $\Omega \subset \R^n$ be a domain with $\Content(\partial \Omega) = 0$. A \textbf{regular partition},
    $\mathcal{P}_{reg} = \{D_{k}\}_{k=1}^{N} $ is a collection of subdomains $D_k \subset \Omega$; $k = 1, \dots, N$ s.t.
    \begin{enumerate}
        \item $\Omega = \bigcup_{k = 1}^N D_k$, $\Content(\partial D_k) = 0 \quad \forall k = 1, \dots, N$.
        \item $\Content(D_k \cap D_\ell) = 0 \quad \forall k \neq \ell$.   
    \end{enumerate}     
\end{definition}

\begin{definition}
    Suppose $f : \Omega \to \R$ is bounded, where $\Omega \subset \R^n$ is a bounded domain, and $\mathcal{P}_{reg} = \{D_k\}_{k = 1}^N$.
    Define 
    \[
        U(f, \mathcal{P}_{reg}) = \sum_{k = 1}^N (\sup_{D_k} f) \cdot \Vol(D_k) \quad L(f, \mathcal{P}_{reg}) = \sum_{k = 1}^N (\inf_{D_k} f) \cdot \Vol(D_k)
    \]    
\end{definition}

\begin{lemma}
    $f : B \subset \R^n \to \R$ is Riemann integrable over box $B$ iff $\forall \epsilon > 0$, $\exists \mathcal{P}_{reg}$ regular partition of $B$ s.t.
    \[
        U(f, \mathcal{P}_{reg}) - L(f, \mathcal{P}_{reg}) \leq \epsilon
    \]   
\end{lemma}

\begin{theorem}[Change of Variables]
    Let $\Omega, \Omega' \subset \R^n$ be bounded domains. Let $T : \Omega' \to \Omega$ be a $C^1$ map that is bijective from $\Omega' $ onto $\Omega$.
    Then, for any $f : \Omega \to \R$ Riemann integrable over $\Omega$,
    \[
        \int_{\Omega} f \; dV = \int_{\Omega'} (f \circ T) \cdot |\det DT| \; dV
    \]       
\end{theorem}

\begin{exmp}[Polar Coordinates]
    Let $x = r \cos \theta$ and $y = r \sin \theta$. $\Omega = \{ (x, y) \in \R^2 : x^2 + y^2 \leq a^2 \}$, $\Omega' = \{(r, \theta) \in (0, a] \times [0, 2 \pi)\}$
    Then $T : \Omega' \to \Omega$ defined by $T(r, \theta) = (r \cos \theta, r \sin \theta)$ is a $C^1$ bijection. Also, $DT = \begin{bmatrix}
        \cos \theta & -r \sin \theta \\
        \sin \theta & r \cos \theta
    \end{bmatrix} $ and $\det DT = r$. Thus, for any $f : \Omega \to \R$ Riemann integrable over $\Omega$,
    \[
        \iint_{\{x^2 + y^2 \leq a^2\}} f(x,y) \; dA = \int_{0}^{a} \int_{0}^{2\pi} f(r \cos \theta, r \sin \theta) \cdot r \; d\theta \; dr
    \]   
\end{exmp}
\begin{remark}
    Change of variables still holds if we only assume $T : \Omega' \setminus \Gamma' \to \Omega$ is a $C^1$ bijection, where $\Content(\Gamma') = 0$  
\end{remark}

\begin{exmp}
    Compute $\iint_{\Omega} (y-x) \cos xy \; dA$ where $\Omega$ is the region defined by $x + y = 1, x + y = 4$, $xy = -1, xy =-4$, $y \geq x$.
    
    To solve, first note that $f(x,y) = (y-x) \cos xy$ is $C^0$ on $\Omega$, thus it is Riemann integrable over $\Omega$.
    Now to solve, we will perform a change of variables with $u = x + y$ and $v = xy$. Then
    $T : \Omega' \to \Omega$ defined by $T(u, v) = (x, y)$ is a $C^1$ bijection, where $\Omega' = \{(u,v) : 1 \leq u \leq 4, -4 \leq v \leq -1\}$.
    Also, $|\det (DT^{-1})| |\det \begin{bmatrix}
        1 & 1 \\
        y & x
    \end{bmatrix}| = |x - y| = \sqrt{u^2 - 4v} = \frac{1}{|\det DT|}$. Thus by change of variables,
    \begin{align*}
        \iint_{\Omega} (y-x) \cos xy \; dA &= \int_{1}^{4} \int_{-4}^{-1} \sqrt{u^2 - 4v} \cos v \cdot \frac{1}{\sqrt{u^2 - 4v}} \; dv \; du \\
        &= \int_{1}^{4} \int_{-4}^{-1} \cos v \; dv \; du = \int_{1}^{4} \left[ \sin v \right]_{-4}^{-1} du = 3(\sin (-1) - \sin(-4))
    \end{align*}
\end{exmp}

\begin{exmp}
    Compute $\iint_D e^{-(x^2 + y^2)} \; dA$ where $D = \{x^2 + y^2 \leq R^2\}$.
    
    To solve, first note that $f(x,y) = e^{-(x^2 + y^2)}$ is $C^0$ on $\Omega$, thus it is Riemann integrable over $\Omega$.
    Now we perform a change of variables with $x = r \cos \theta$ and $y = r \sin \theta$. Then $T : \Omega' \to \Omega$ is defined as
    $T(r, \theta) = (r \cos \theta, r \sin \theta)$. Then $T$ is a $C^1$ bijection, where $\Omega' = \{(r, \theta) : 0 \leq r \leq R, 0 \leq \theta \leq 2\pi\}$.
    Also, $DT = \begin{bmatrix}
        \cos \theta & -r \sin \theta \\
        \sin \theta & r \cos \theta
    \end{bmatrix} $ and $\det DT = r$. Thus by change of variables,
    \[
        \iint_{D} e^{-(x^2 + y^2)} \; dA = \int_{0}^{R} \int_{0}^{2\pi} e^{-r^2} \cdot r \; d\theta \; dr = 2\pi \int_0^R r e^{-r^2} \; dr = \pi (1 - e^{-R^2})
    \]    
\end{exmp}

\begin{exmp}[spherical coordinates]
    Let $x = p \sin \varphi \cos \theta$, $y = p \sin \varphi \cos \theta$, $z = p \cos \theta$
    Let $\Omega = \{(x, y, z) : x^2 + y^2 + z^2 \leq R^2\}$ and $\Omega' = \{(p, \varphi, \theta) : 0 \leq p \leq a, 0 \leq \varphi \leq \pi, 0 \leq \theta \leq 2 \pi\}$.
    Then for any $f : \Omega \to \R$ Riemann integrable over $\Omega$,
    \[
        \iiint_{\Omega} f \; dV = \iint_{\Omega'} f(p \sin \varphi \cos \theta, p \sin \varphi \sin \theta, p \cos \varphi) \cdot p^2 \sin \varphi \; dp \; d\theta \; d\varphi
    \] 
\end{exmp}

\subsection{Improper Integrals}


\begin{definition}
    Suppose $A \subset \R^n$ a region, $f : A \to \R$, $f(x) \geq 0 \quad \forall x \in A$.
    Suppose $\forall r > 0$, $f$ is bounded and integrable over $A \cap B_r(0)$. Then we define
    \[
        \int_{A} f \; dV := \lim_{r \to \infty} \int_{B_r(0) \cap A} f \; dV
    \]      
    If it exists, is is the \textbf{improper integral} of $f$ over $A$.
\end{definition}

\begin{definition}
    Consider $\{A_k\}_{k = 1}^\infty$, $A_k \subset A$, $A_k \uparrow$. Then $\{A_k\}_{k =1}^\infty$ are
    \textbf{exhausting regions} of $A$ if 
    \begin{enumerate}
        \item $A_k$ compact for all $k$;
        \item $\Content(\partial A_k) = 0$;
        \item $\forall r > 0$, $\exists k_0(r)$ s.t. $B_r(0) \cap A \subset A_k$, $\forall k \geq k_0(r)$.     
    \end{enumerate} 

    For example, $\{B_k(0) \cap A\}_{k = 1}^\infty$ are exhausting regions of $A$.
\end{definition}

\begin{theorem}
    Assume $f \geq 0$ bounded and integrable on $B_r(0) \cap A$, $\forall r > 0$. Then for any exhausting regions
    $\{A_k\}$ of $A$, 
    \[
        \lim_{k \to \infty} \int_{A_k} f \; dV = \lim_{r \to \infty} \int_{B_r(0) \cap A} f \; dV
    \]      
\end{theorem}

\begin{exmp}
    Evaluate $\int_0^\infty e^{-x^2}$.
    
    To solve, first note that $f(x) = e^{-x^2}$ is bounded and integrable over $[0, r]$ for all $r > 0$. Thus $\int_0^\infty e^{-x^2}$ exists. 
    Also, $\{[0, k]\}_{k = 1}^\infty$ are exhausting regions of $[0, \infty)$. Thus,
    \begin{align*}
        \left(\int_0^\infty e^{-x^2} \; dx \right)^2 = \int_0^\infty e^{-x^2} \; dx \cdot \int_0^\infty e^{-y^2} \; dy &= \lim_{k \to \infty} \iint_{[0, k]^2} e^{-(x^2 + y^2)} \; dA \\
        &= \lim_{k \to \infty} \int_0^k \int_0^k e^{-(x^2 + y^2)} \; dy \; dx = \lim_{k \to \infty} \frac{\pi}{2} \int_0^k r e^{-r^2} \; dr = \frac{\pi}{4}
    \end{align*}
    Therefore $\int_0^\infty e^{-x^2} \; dx = \frac{\sqrt{\pi}}{2}$.
\end{exmp}

\begin{definition}
    Given $f : A \to \R$, define $f^+ = \max\{f, 0\}$ and $f^- = \max\{-f, 0\}$. Then $f = f^+ - f^-$ and $|f| = f^+ + f^-$. 
    iF both $\int_A f^+ \; dV$ and $\int_A f^- \; dV$ exist, then f integrable over $A$ and 
    \[
        \int_{A} f \; dV = \int_A f^+ \; dV - \int_A f^- \; dV
    \] 
\end{definition}

\section{Vector Fields, Curves, and Surfaces}

\subsection{Path Integrals/ Line Integrals}

\begin{definition}
    A mapping $\gamma : [a,b] \to \R^n$ is called $C^1([a,b])$ if $\gamma \in C^1(a,b)$ and
    \[
        \lim_{t \to a^+} \frac{\gamma(t) - \gamma(a)}{t - a} = \gamma'(a^+) \quad \text{and} \quad \lim_{t \to b^-} \frac{\gamma(t) - \gamma(b)}{t - b} = \gamma'(b^-)
    \]  

    $\gamma : [a,b] \to \R^n$ is a \textbf{parameteric curve} in $\R^n$ if $\gamma \in C^1([a,b])$.
    Then
    \[
        \gamma(t) = (x_1(t), x_2(t), \dots, x_n(t)) \quad \text{ and } \quad \gamma'(t) = (x_1'(t), x_2'(t), \dots, x_n'(t)) \text{ is the tangent vector}
    \] 
\end{definition}

\begin{definition}
    Let $\gamma : [a,b] \to \R^n$ be a continuous map. Then $\gamma$ is \text{piecewise} $C^1$ if it is $C^1$ on each subinterval of a partition of $[a,b]$.
    i.e. it has finitely many "corners"/non-differentiable points.

    We say $\gamma $ is \textbf{regular} if $\gamma \in C^1([a,b])$ and $\gamma'(t) \neq 0$
    for all $t \in [a,b]$. We say $\gamma$ is \textbf{piecewise regular} if it is regular on each subinterval of a partition of $[a,b]$.    
\end{definition}

\begin{prop}
    Let $\gamma :[a,b] \to \R^n$ be regular. Then there exists a partition $a = t_0 < t_1 < \dots < t_N = b$
    s.t. on each $[t_{i -1}, t_i]$, $i = 1, \dots, N$, there is a $j$ s.t. $t \mapsto x_{j_i}(t)$ is $C^1$ invertible
    and for $t \in [t_{i -1}, t_i]$, $\gamma(t) = (x_1(x_{j_i}(t)), \dots, x_{j_{i - 1}}(x_{j_i}(t)), x_{j_i}(t), x_{j_{i + 1}}(x_{j_i}(t)),
     \dots, x_n(x_{j_i}(t)))$. i.e. $\gamma$ is a graph over $x_{j_i}$.           
\end{prop}

\begin{definition}
    For each piecewise $C^1$ curve $\gamma$, we define the \textbf{length} of $\gamma$ as
    \[
        L(\gamma) := \sum_{i = 1}^n \int_{a_{i-1}}^{a_i} \| \gamma'(t) \| \; dt = \sum_{i =1}^n \int_{a_{i-1}}^{a_i} \sqrt{(x_1'(t))^2 + \dots + (x_n'(t))^2} \; dt \quad \footnote{Proof follows from MVT} 
    \]   
\end{definition}

\begin{definition}[Reparameterization]
    Let $\gamma : [a,b] \to \R^n$ be $C^1$. $\tilde{\gamma} : [c,d] \to \R^n$ is a \textbf{reparameterization} of $\gamma$
    if there exists a $C^1$ bijection $h : [c,d] \to [a,b]$ s.t. $\tilde{\gamma}(t) = \gamma(h(t))$.
    
    Note then $\tilde{\gamma}'(t) = h'(t) \cdot \gamma'(h(t))$. Thus, if $\gamma$ is regular and $h'(t) > 0$ for all $t$, then $\tilde{\gamma}$ is also regular.
\end{definition}

\begin{theorem}
    Suppose $\gamma :[a,b] \to \R^n$ is a $C^1$ curve, and $\tilde{\gamma} : [c,d] \to \R^n$ is a $C^1$ reparameterization.
    Then $L(\gamma) = L(\tilde{\gamma})$.    
\end{theorem}

\begin{definition}[arc-length parameterization]
    Let $\gamma : [a,b] \to \R^n$ be a $C^1$ regular parameterized curve. Then define
    \[
        s(t) = \int_a^t \| \gamma'(u) \| \; du \in C^1([a,b]) \quad \text{and} \quad s'(t) = \| \gamma'(t) \| > 0 \quad \forall t \in [a,b] 
    \]   
    Then $s : [a,b] \to [0, L]$ is a $C^1$ monotone increasing bijection where $L = L(\gamma) = \int_a^b \| \gamma'(t) \| \; dt$. Thus, $s$ has a $C^1$ inverse $t(s) : [0, L] \to [a,b]$. Then $\tilde{\gamma}(s) = \gamma(t(s))$ is a reparameterization of $\gamma$ and is called the \textbf{arc-length parameterization} of $\gamma$. Note that $\|\tilde{\gamma}'(s)\| = 1$ for all $s$.  
\end{definition}

\begin{definition}
    Let $f$ be $C^0$ in a neighbourhood of $\gamma$. The \textbf{path integral} of $f$ along $\gamma$ is defined as
    \[
        \int_{\gamma} f \; ds := \int_a^b f(\gamma(t)) \cdot \| \gamma'(t) \| \; dt
    \]     
    If we instead use arc-length to parameterize the curve, then
    \[
        \int_{\gamma} f \; ds = \int_0^L f(\gamma(t(s))) \; ds 
    \] 
    by the change of variables formula, since $s'(t) = \| \gamma'(t) \|$ 
\end{definition}

\subsection{Vector Fields and Line Integrals}

\begin{definition}[Vector Field]
    A \textbf{vector field} on $\Omega \subset \R^n$ is a map $\vec{F} : \Omega \to \R^n$. 
\end{definition}

\begin{definition}
    The \textbf{total work} done by $\vec{F}$ along $\gamma$ is defined as
    \[
        \int_{\gamma} \vec{F} \; ds := \lim_{N \to \infty} \sum_{i = 1}^N \vec{F}(\gamma(t_i)) \; \Delta s_i 
        \quad \text{ and } \quad = \int_a^b \vec{F}(\gamma(t)) \cdot \gamma'(t) \; dt \quad \text{assuming the integral exists.}
    \] 
\end{definition}

\begin{definition}[Line Integral]
    Let $\gamma \subset \R^n$ be a piecewise smooth path and let $\vec{F}$ be a vector field defined on a neighbourhood of $\gamma$.
    The \textbf{line integral} of $\vec{F}$ along $\gamma$ is defined as
    \[
        \int_{\gamma} \vec{F} \cdot \vec{ds} = \int_a^b \vec{F}(\gamma(t)) \cdot \gamma'(t) \; dt \quad \text{if } \gamma \text{ is regular.}
    \]     
    Note if we let $\vec{T}(t) = \frac{\gamma'(t)}{\| \gamma'(t) \| }$ be the unit tangent field, then 
    \[
        \int_{\gamma} \vec{F} \cdot \vec{ds} = \int_{\gamma} \vec{F}(\gamma(t)) \cdot \vec{T}(t) \; ds \quad \text{where } ds = \| \gamma'(t) \| dt
    \]  
\end{definition}

\begin{definition}
    Let $\gamma :[a,b] \to \R^n$ be a piecewise $C^1$ path and $\tilde{\gamma} : [c,d] \to \R^n$ a reparameterization of
    $\gamma$ with a bijective $C^1$ map $h : [c,d] \to [a,b]$ s.t. $\tilde{\gamma}(t) = \gamma(h(t))$. Then
    \begin{enumerate}
        \item $\tilde{\gamma}$ is \textbf{orientation-preserving} if $\tilde{\gamma}(c) = \gamma(a)$ and $\tilde{\gamma}(d) = \gamma(b)$;
        \item $\tilde{\gamma}$ is \textbf{orientation-reversing} if $\tilde{\gamma}(c) = \gamma(b)$ and $\tilde{\gamma}(d) = \gamma(a)$.
    \end{enumerate}     
\end{definition}

\begin{lemma}
    Let $\gamma : [a,b] \to \R^n$ be a piecewise smooth path and $\tilde{\gamma} : [c,d] \to \R^n$ a reparameterization 
    of $\gamma$. Let $\vec{F}$ be a $C^0$ vector field. Then
    \[
        \int_{\gamma} \vec{F} \cdot \vec{ds} = \epsilon \int_{\tilde{\gamma}} \vec{F} \cdot \vec{ds} \quad \text{where } \epsilon = 1 \text{ if } \tilde{\gamma} \text{ is orientation-preserving and } \epsilon = -1 \text{ otherwise.}
    \]    
\end{lemma}

\subsection{Conservative Vector Fields}

\begin{definition}[Conservative Vector Field]
    $\vec{F} \in C^1(\Omega, \R^n)$ is \textbf{conservative} if $\exists f \in C^1(\Omega)$ s.t.
    $\vec{F} = \nabla f$. 
\end{definition}

\begin{exmp}
    $\vec{F}(x,y) = (x,y)$, $(x, y) \in \R^2$ is the Radial Vector Field. Here, set $f(x,y) = \frac{1}{2}(x^2 + y^2)$.
    Then $\nabla f = (x, y) = \vec{F}$, so $\vec{F}$ is conservative.       
\end{exmp}

\begin{theorem}
    Suppose $\vec{F}$ is conservative in a neighbourhood of a piecewise smooth path $\gamma$.
    Then 
    \[
        \int_{\gamma} \vec{F} \cdot \vec{ds} = f(\gamma(b)) - f(\gamma(a))
    \]   
\end{theorem}
\begin{corollary}
    If $\vec{F}$ is conservative in a neighbourhood of a closed piecewise smooth path $\gamma$, then $\int_{\gamma} \vec{F} \cdot \vec{ds} = 0$.
\end{corollary}

\begin{theorem}
    Assume $\vec{F} = (F_1, F_2, \ldots, F_n) \in C^1(\Omega, \R^n)$ is conservative. Then
    \[
        \frac{\partial F_i}{\partial x_j}(x) = \frac{\partial F_j}{\partial x_i}(x) \quad \forall i, j = 1, \dots, n, \forall x \in \Omega
    \]   
\end{theorem}

\subsection{Poincaré's Lemma}

\begin{lemma}
    Suppose $g(x,y)$ is continuous defined in $\Omega \times A$, where $\Omega \subset \R^n$ open and
    $A \subset \R^m$ compact. Assume in addition that $g \in C^1(\Omega \times A)$. Then
    \[
        \frac{\partial}{\partial x_i}\left(\int_{A} g(x,y) \; dy \right) = \int_A \frac{\partial }{\partial x_i}g(x,y) \; dy \quad \forall i = 1, \dots, n
    \]      
\end{lemma}

\begin{definition}
    $\Omega \subset \R^n$ is \textbf{star-shaped} if $\exists p_0 \in \Omega$ s.t. the line $\ell(p_0, x)$ is
    contained in $\Omega$ for all $x \in \Omega$. i.e. $\forall x \in \Omega$, $\{(1 - t)p_0 + tx : t \in [0, 1]\} \subset \Omega$.   
\end{definition}

\begin{theorem}[Poincaré's Lemma]
    Suppose $\Omega \subset \R^n$ is star shaped and $\vec{F} \in C^1(\Omega, \R^n)$. Then $\vec{F}$ is conservative
    iff 
    \[
        \frac{\partial F_i}{\partial x_j}(x) = \frac{\partial F_j}{\partial x_i}(x) \quad \forall i, j = 1, \dots, n, \forall x \in \Omega
    \]    
\end{theorem}


\subsection{Surfaces and Surface Integrals}


\begin{definition}[Parameterized Surface]
    A \textbf{parameterized surface} $S \subset \R^3$ is the image of a map $\Phi : D \to \R^3 $,
    \[
        S = \left\{ (x(u,v), y(u,v), z(u,v)) = \Phi(u,v) : (u,v) \in D \right\} 
    \]   
    Where $D \subset \R^2$ is some bounded domain. If $\Phi$ is $C^k$, we call $S$ a $C^k$ surface. Note that:
    \[
        \Phi_u = (\frac{\partial x}{\partial u}, \frac{\partial y}{\partial u}, \frac{\partial z}{\partial u}) \quad \text{and} \quad \Phi_v = (\frac{\partial x}{\partial v}, \frac{\partial y}{\partial v}, \frac{\partial z}{\partial v})
    \]     
    are the tangent vectors of $S$ at $\Phi(u,v)$. The tangent plane of $S$ at $\Phi(u,v)$ is the plane spanned by $\Phi_u$ and $\Phi_v$.
\end{definition}

\begin{definition}[Regular Surface]
    We say $S$ is a \textbf{regular surface} if $\vec{\Phi_u} \times \vec{\Phi_v} \neq \vec{0}$ for all $(u,v) \in D$. 

\end{definition}

\begin{definition}[Area of a Surface]
    Let $S$ be a regular $C^1$ surface parameterized by $\Phi : D \to S$. Then the area of $S$ is defined as
    \[
        \operatorname{area}(S) := \iint_{D} \| \vec{\Phi_u} \times \vec{\Phi_v} \| \; dA
    \] 
    Here,
    \[
        \| \vec{\Phi_u} \times \vec{\Phi_v} \| = \sqrt{\left| \frac{\partial(x,y)}{\partial(u,v)} \right|^2 + \left| \frac{\partial(y,z)}{\partial(u,v)} \right|^2 + \left| \frac{\partial(x,z)}{\partial(u,v)} \right|^2}
    \] 
    
\end{definition}

\begin{definition}
    Let $S = \Phi(D)$ be a $C^1$ parameterized surface and suppose $f \in C^0(S)$, i.e. $f$ is $C^0$ in a neighbourhood of $S$ in $\R^3$, i.e. $f \circ \Phi \in C^0(D, \R^3)$.
    Then the \textbf{surface integral} of $f$ over $S$ is defined as
    \[
        \iint_S f \; dS := \iint_{D} f(\Phi(u,v)) \cdot \| \vec{\Phi_u} \times \vec{\Phi_v} \| \; du \; dv
    \]        
\end{definition}

\begin{exmp}
    There is an important special case. Suppose $S = \left\{ (x, y, g(x,y)) : (x,y) \in D \text{ and } g \in C^1(D) \right\} $.
    Here $(x,y) = (u,v)$, $\Phi(u,v) = (u, v, g(u,v))$, $\Phi \in C^1(D, \R^3)$. In this case
    \[
        \vec{\Phi_x} = (1, 0, g_x) \quad \text{and} \quad \vec{\Phi_y} = (0, 1, g_y) \quad \text{and} \quad \| \vec{\Phi_x} \times \vec{\Phi}_y \| = \sqrt{1 + g_x^2 + g_y^2} 
    \]    
    Then $dS = \sqrt{1 + g_x^2 + g_y^2} \; dx \; dy$ and so
    \[
        \iint_S f \; dS = \iint_D f(x, y, g(x,y)) \sqrt{1 + g_x^2 + g_y^2} \; dx \; dy
    \] 
\end{exmp}

\begin{exmp}[Helicoid]
    $D = \{0 \leq \theta \leq 2 \pi , 0 \leq r \leq 1\}$. $\Phi(r, \theta) = (r \cos \theta, r \sin \theta, \theta)$ is a $C^1$ parameterization of the helicoid. Then
    \[
        \vec{\Phi_r} = (\cos \theta, \sin \theta, 0) \quad \text{and} \quad \vec{\Phi_\theta} = (-r \sin \theta, r \cos \theta, 1) \quad \text{and} \quad \| \vec{\Phi_r} \times \vec{\Phi_\theta} \| = \sqrt{1 + r^2}
    \] 
\end{exmp}

\begin{definition}
    Let $S = \Phi(D)$ be a $C^1$ parameterized surface. We say that $S$ is \textbf{regular} if $(\vec{\Phi_u} \times \vec{\Phi_v})(u,v) \neq \vec{0}$ for all $(u,v) \in D$.
    If $S$ is regular, then we can define the \textbf{unit normal} 
    \[
        \vec{n} = \frac{\pm (\Phi_u \times \Phi_v)}{\|\Phi_u \times \Phi_v\|}
    \]  
    The choice of $\pm$ is called the \textbf{orientation} of $S$. Note that if $S$ is regular, then $\vec{n}$ is well-defined and continuous on $S$.
\end{definition}

\begin{exmp}
    $S$ is the graph of $g$, i.e. $S = \{(x, y, g(x,y)): (x,y) \in D \text{ and } g \in C^1(D) \}$. Then
    \[
        \vec{n} = \frac{\pm (-g_x, -g_y, 1)}{\sqrt{1 + g_x^2 + g_y^2}}
    \]
\end{exmp}

\begin{definition}[Surface Integral of a Vector Field]
    Let $F$ be a $C^0$ vector field on an oriented regular parameterized surface. Then
    \[
        \iint_S \vec{F} \cdot \vec{dS} := \iint_D \vec{F}(\Phi(u,v)) \cdot \vec{n} \; ds = \iint_D \vec{F}(\Phi(u,v)) \frac{\Phi_u \times \Phi_v}{\|\Phi_u \times \Phi_v\|} \cdot \| \Phi_u \times \Phi_v \|  \; du \; dv = \iint_D \vec{F}(\Phi(u,v)) \cdot (\Phi_u \times \Phi_v) \; du \; dv
    \] 
\end{definition}

\begin{remark}
    The definition is independent of the choice of parameterization, modulo the choice of oerientation.
\end{remark}

\begin{definition}[reparameterization of a surface]
    Let $\Phi : D \subset \R^2 \to S$ be a regular $C^1$ parameterization, $\tilde{\Phi} : \tilde{D} \subset \R^2 \to S$ another regular $C^1$ parameterization.
    We say $\tilde{\Phi}$ is a \textbf{reparameterization} of $\Phi$ if there exists a $C^1$ bijection $h : \tilde{D} \to D$ s.t. $\tilde{\Phi}(u,v) = \Phi(h(u,v))$. Then
    $\tilde{\Phi}$ is \textbf{orientation-preserving} (+) if 
    \[
        \frac{\Phi_u \times \Phi_v (h(s,t))}{\| \Phi_u \times \Phi_v \| \cdot (h(s,t)) } = \frac{\tilde{\Phi}_s \times \tilde{\Phi}_t(s,t)}{\| \tilde{\Phi}_s \times \tilde{\Phi}_t \| (s,t)} \quad \forall (s,t) \in \tilde{D}
    \] 
    and \textbf{orientation-reversing} (-) if 
    \[
        \frac{\Phi_u \times \Phi_v (h(s,t))}{\| \Phi_u \times \Phi_v \| \cdot (h(s,t)) } = -\frac{\tilde{\Phi}_s \times \tilde{\Phi}_t(s,t)}{\| \tilde{\Phi}_s \times \tilde{\Phi}_t \| (s,t)} \quad \forall (s,t) \in \tilde{D}
    \] 
    
\end{definition}

\begin{lemma}
    Let $\tilde{\Phi}$ and $\Phi$ be as above with $h : D \to \tilde{D}$ such that 
    $h(s,t) = (u(s,t), v(s.t)) \in C^1(\tilde{D}, D)$. Then $h$ is orientation-preserving
    iff $\det (\frac{\partial (u,v)}{\partial (s,t)}) > 0$, and $h$ is orientation-reversing iff
    $\det (\frac{\partial (u,v)}{\partial (s,t)}) < 0$.
\end{lemma}

\begin{prop}
    Let $S$ be an oriented $C^1$ regular parameterized surface $S = \Phi(D)$ and 
    $\vec{F} \circ \Phi \in C^0(D, \R^3)$. Then $\iint_S \vec{F} \cdot \vec{dS}$ is well defined
    independent of the choice of parameterization, modulo the choice of orientation, (+) or (-).    
\end{prop}

\begin{corollary}
    $S = \Phi(D)$ an oriented regular parameterized surface. Then 
    \[
        \iint_{S^\mp} \vec{F} \cdot \vec{dS} = - \iint_{S^\pm} \vec{F} \cdot \vec{dS}
    \]  
\end{corollary}

\subsection{Green's Theorem}

\begin{definition}[Regular Surface]
    $S \subset \R^3$ is a \textbf{regular surface} with $C^1$ boundary if $S \subset \R^3$ is
    compact and if $\forall p \in S$, $\exists r, \delta > 0$, $B_r(p) \subset \R^3$, and
    $D = \{ -\delta < u < \delta, -\delta < v < \delta\} \subset \R^3$ and a regular map
    $\Phi : D \to \R^3$ such that either
    \begin{enumerate}
        \item $S \cap B_r(p) = \Phi(D)$;
        \item $S \cap B_r(p) = \Phi(D^+)$, where $D^+ = \{ -\delta < u < \delta, 0 \leq v < \delta\}$   
    \end{enumerate}       
    Points in (1) are called \textbf{interior points} of $S$ and points in (2) are called \textbf{boundary points} of $S$.
    $\partial S := \{\text{all boundary points of } S\}$    
\end{definition}

\begin{exmp}[orientability]
    Let $S_+^2 = \{x^2 + y^2 + z^2 = 1, z \geq 0\}$. We say that a regular surface $S$ is \textbf{orientable}
    if there exists a unit normal $\vec{n}(p)$ at each $p \in S$ with $\vec{n} \in C^0(S, \R^3)$.     
\end{exmp}

\begin{definition}
    We say a bounded planar domain $D \subset \R^2$ is (+) \textbf{positively oriented} if
    $\partial D$ is oriented counter-clockwise, i.e. if the unit normal is $\vec{k} = (0,0,1)$. 
\end{definition}

\begin{theorem}[Green's Theorem]
    Let $D \subseteq \R^2$ be a domain with piecewise $C^1$ boundary $\partial D$, and
    $\vec{F} = (P, Q) : D \to \R^2$ a $C^1$ vector field on $D$. Then
    \[
        \int_{\partial D} \vec{F} \cdot \vec{ds} = \iint_D \left( \frac{\partial Q}{\partial x} - \frac{\partial P}{\partial y} \right) \; dA
    \]      
\end{theorem}

\begin{corollary}
    \[
        \int_{\partial D} x \; dy = - \int_{\partial D} y \; dx = \frac{1}{2}\int_{\partial D} x \; dy - y \; dx = \operatorname{area}(D)
    \] 
\end{corollary}

\begin{exmp}
    Compute the area of the ellipse $\left\{ \frac{x^2}{a^2} + \frac{y^2}{b^2} \leq 1 \right\} $ using Green.

    To solve, notice the boundary $\partial D$ is the ellipse $\left\{ \frac{x^2}{a^2} + \frac{y^2}{b^2} = 1 \right\}$. We can parameterize $\partial D$ by $\gamma(t) = (a \cos t, b \sin t)$, $t \in [0, 2\pi]$. Then
    \[
        \int_{\partial D} x \; dy = \int_{0}^{2\pi} a \cos t \cdot b \cos t \; dt = ab \int_0^{2\pi} \cos^2 t \; dt = ab \pi
    \]  
\end{exmp}

\begin{corollary}
    Let $D = D_1 \setminus D_2$ and $\partial D_1$, $\partial D_2$ piecewise $C^1$ with orientations
    induced from $D_1 + D_2$. Then if $\vec{F} = (P, Q)$ is a $C^1$ vector field on $D$ with 
    $\frac{\partial Q}{\partial x} - \frac{\partial P}{\partial y} = 0$ on $D$, then
    \[
        \int_{\partial D_1} \vec{F} \cdot \vec{ds} = \int_{\partial D_2} \vec{F} \cdot \vec{ds}
    \] 
\end{corollary}

\begin{theorem}
    If $\vec{F} = (P, Q) \in C^1$, then 
    \[
        \int_{\partial D} \vec{F} \cdot \vec{n} \; ds = \iint_D \left( \frac{\partial P}{\partial x} + \frac{\partial Q}{\partial y} \right) \; dA
    \]  
\end{theorem}

\subsection{Harmonic Functions}

\begin{definition}[Harmonic Function]
    Let $\Delta = \frac{\partial^2}{\partial x^2} + \frac{\partial^2}{\partial y^2}$ be the Laplacian. $u \in C^2(D)$ is 
    \textbf{harmonic} if $\Delta u = 0$ on $D$.  
\end{definition}

\begin{corollary}
    $u$ harmonic on domain $D \subset \R^2$ with piecewise $C^1$ boundary $\partial D$. Then  
    \[
        \iint_{D} \Delta u \; dA = \int_{\partial D} \Delta u \cdot \vec{n} \; ds = 0
    \]  
\end{corollary}

\begin{corollary}[Green's Identities]
    Let $\varphi, \psi \in C^2(D)$, $\partial D$ piecewise $C^1$ with positive orientation. Then
    \begin{itemize}
        \item \[ \iint_D \varphi \Delta \psi \; dA = - \int_{\partial D} \nabla \varphi \cdot \nabla \psi \; dA + \iint_{\partial D} \varphi \nabla \psi \cdot \vec{n} \; dA \]
        \item \[ \iint_D \left[ \varphi \Delta \psi- \psi \Delta \varphi \right] \; dA = \int_{\partial D} \left( \varphi \nabla \psi - \psi \nabla \varphi \cdot \vec{n} \right) ds \]
    \end{itemize}    
    
\end{corollary}

\subsection{Stokes' Theorem}


\begin{definition}
    Let $\vec{F} = (F_1, F_2, F_3): \Omega \subset \R^3 \to \R^3$ be a $C^1$ vector field. Then the \textbf{curl} of $\vec{F}$ is defined as
    \[
        \curl \vec{F} = \nabla \times \vec{F} = \det \begin{bmatrix}
        \vec{i} & \vec{j} & \vec{k} \\
        \frac{\partial}{\partial x_1} & \frac{\partial}{\partial x_2} & \frac{\partial}{\partial x_3} \\
        F_1 & F_2 & F_3
        \end{bmatrix}
        = \left( \frac{\partial F_3}{\partial x_2} - \frac{\partial F_2}{\partial x_3}, \frac{\partial F_1}{\partial x_3} - \frac{\partial F_3}{\partial x_1}, \frac{\partial F_2}{\partial x_1} - \frac{\partial F_1}{\partial x_2} \right)
    \]    
\end{definition}

\begin{theorem}[Stokes]
    Let $S$ be a compact regular oriented surface with (+) orientation and piecewise $C^1$ boundary $\partial S$. Let $F : S \to \R^3$ be a $C^1$ vector field.
    Then
    \[
        \iint_S \curl \vec{F} \cdot \vec{dS} = \int_{\partial S} \vec{F} \cdot \vec{ds}
    \]      
\end{theorem}

\begin{remark}
    Note, 
    \[
        \iint_S \curl \vec{F} \cdot \vec{dS} = \iint_S \curl \vec{F} \cdot \vec{n} \; dS
    \] 
\end{remark}

\begin{exmp}
    $\vec{F} = (P, Q, 0)$. Then $F_3 = 0$ and $P = P(x, y), Q = Q(x, y)$, so $\curl \vec{F} = \left( 0, 0, \frac{\partial Q}{\partial x} - \frac{\partial P}{\partial y} \right)$.
    Then we can see Green's theorem is a really just a special case of Stokes' theorem, since
    \[
        \left( 0, 0, \frac{\partial Q}{ \partial x} - \frac{\partial P}{\partial y} \right) \cdot \vec{k} = \frac{\partial Q}{\partial x} - \frac{\partial P}{\partial y}
    \]    
\end{exmp}

\begin{corollary}
    Suppose $S_1, S_2$ are two oriented surfaces with $\partial S_1 = \partial S_2 = \gamma$, and suppose $S_1$ and $S_2$ induce 
    opposite orientations on $\gamma$. Then, for any $\vec{F} \in C^1$,
    \[
        \iint_{S_1} \curl \vec{F} \cdot \vec{dS} = - \iint_{S_2} \curl \vec{F} \cdot \vec{dS}
    \]      
\end{corollary}

\begin{corollary}
    If $S$ is a compact, oriented surface with $\partial S = \emptyset$ and $\vec{F} \in C^1(S)$, then
    \[
        \iint_S \curl \vec{F} \cdot \vec{dS} = 0
    \]  
\end{corollary}

\begin{corollary}
    Suppose $\vec{F} = \nabla F$ for some $f \in C^2$ and $S$ is an oriented surface with $\partial S = \gamma_1 - \gamma_2$.
    Then,
    \[
        \int_{\gamma_1} \vec{F} \cdot \vec{ds} = \int_{\gamma_2} \vec{F} \cdot \vec{ds}
    \]     
\end{corollary}

\begin{exmp}
    Let $C = \{ y + z = 2\} \cap \{x^2 + y^2 = 1 \}$ with (+) orientation. Let $\vec{F} = (-y^2, x, z^2)$. Evaluate $\int_C \vec{F} \cdot \vec{ds}$.
    
    To solve, by Stokes
    \[
        \int_C \vec{F} \cdot \vec{ds} = \iint_S \curl \vec{F} \cdot \vec{dS}
    \] 
    Here $\curl \vec{F} = (0, 0, 1 + 2y)$ and $\vec{n} = \frac{(0, 1, 1)}{\sqrt{2}}$ 
    Therefore,
    \[
        \int_C \vec{F} \cdot \vec{ds} = \iint_D (0, 0, 1 + 2y) \cdot (0, 1, 1) \; dA = \iint_D (1 + 2y) \; dA
    \] 
    To solve this integral, we make the change of variables $x = r \cos \theta, y = r \sin \theta$, $D = \{0 \leq r \leq 1, 0 \leq \theta \leq 2 \pi \}$.
    Then we get
    \[
        = \int_0^{2 \pi} \int_0^1 (1 + 2r \sin \theta) r \; dr \; d \theta = \int_0^{2 \pi} \left( \frac{1}{2} + \frac{2}{3} \sin \theta \right) d\theta = \pi
    \]   
\end{exmp}

\subsection{Divergence Theorem}

\begin{theorem}[Divergence Theorem]
    Let $W \subset \R^3$ be a bounded solid region in $\R^3$ with piecewise smooth boundary $S = \partial W$ with positive orientation and
    $\vec{F} \in C^1(W, \R^3)$. Then
    \[
        \iiint_W \dive \vec{F} \; dV = \iint_S \vec{F} \cdot \vec{dS}
    \]     
    Here, $\dive \vec{F} := \frac{\partial F_1}{\partial x} + \frac{\partial F_2}{\partial y} + \frac{\partial F_3}{\partial z}$
\end{theorem}

\begin{corollary}
    Under the same assumptions as above, for any $f, g \in C^2(W)$, 
    \[
        \iiint_W \left[g \Delta f - f \Delta g \right] \; dV = \iint_{\partial W} \left[ g \nabla f \cdot \vec{dS} - f \nabla g \cdot \vec{dS} \right]
    \]  
\end{corollary}

\begin{corollary}
    Suppose $W, W_1, W_2, \dots, W_k \subset \R^n$ are regions as in the divergence theorem with
    \begin{itemize}
        \item $W_i \cap W_j = \emptyset$ for all $i \neq j$;
        \item $W_j \subset W^{\text{int}}$, $p_j \in W_j^{\text{int}}$ for all $j = 1, \dots, k$;  
    \end{itemize}
    for some points $p_1, \dots, p_k$. If $\vec{F} \in C^1(W \setminus \{p_1, \dots, p_k\})$ and $\dive \vec{F} = 0$ for all $x \in W \setminus \{p_1, \dots, p_k\}$, then
    \[
        \iint_{\partial W} \vec{F} \cdot \vec{dS} = \sum_{j = 1}^k \iint_{\partial W_j} \vec{F} \cdot \vec{dS}
    \]  
\end{corollary}

\begin{exmp}[Gauss law]
    Suppose $W \subset \R^3$ as in the divergence theorem and $\vec{F} = (x, y, z)$ and $r = \sqrt{x^2 + y^2 + z^2}$. Then, if $(0, 0, 0) \notin \partial W$,
    \[
        \iint_{\partial W} \frac{\vec{F}}{r^3} \cdot \vec{dS} = \begin{cases}
        4 \pi & \text{if } (0, 0, 0) \in W^{\text{int}} \\
        0 & \text{otherwise}
        \end{cases}
    \]    
    To show this, set $\vec{F} = \frac{\vec{r}}{r^3}$. Then $\dive \vec{F} = 0$, $\vec{r} \in \R^3 \setminus (0, 0, 0)$. SO if $(0, 0, 0) \notin W$, $\vec{F} \in C^1(W)$ and by
    the divergence theorem, $\iint_{\partial W} \frac{\vec{r}}{r^3} \cdot \vec{dS} = 0$. If $(0, 0, 0) \in W$, pick $\delta > 0$ small so that $B_{\delta}(0) \subset W^{\text{int}}$.
    Let $\tilde{W} = W \setminus B_{\delta}(0)$. By the most recent corollary,
    \[
        \iint_{\partial W} \frac{\vec{r}}{r^3} \cdot \vec{dS} = \iint_{\partial B_{\delta}(0)} \frac{\vec{r}}{r^3} \cdot \vec{dS} = \iint_{\partial B_{\delta}(0)} \frac{\vec{r}}{\delta^3} \cdot \vec{dS} = \frac{1}{\delta^3} \iint_{\partial B_{\delta}(0)} \vec{r} \cdot \vec{dS} = 4 \pi
    \]                
\end{exmp}

\subsection{Divergence Theorem Applications}

\begin{exmp}
    Compute $\iint_S \vec{F} \cdot \vec{dS}$ where $\vec{F} = (3x + z^{77}, y^2 - sinx^2 z, xz + ye^{x^3})$,  
    $S = \partial W$ where $W = \{0 \leq x \leq 1, 0 \leq y \leq 3, 0 \leq z \leq 2 \}$.
    
    To solve, we can use the divergence theorem. We have $\dive F = 3 + 2y + x$. Then by the divergence theorem,
    \begin{align*}
        \iint_S \vec{F} \cdot \vec{dS} = \iiint_W (3 + 2y + x) \; dV &= \int_0^1 \int_0^3 \int_0^2 (3 + 2y + x) \; dz \; dy \; dx \\
        &= \int_0^1 \int_0^3 2(3 + 2y + x) \; dy \; dx \\
        &= \int_0^1 2(9 + 9 + 3x) \; dx \\
        &= \int_0^1 36 + 6x \; dx = 39
    \end{align*} 
\end{exmp}

\end{document}